\chapter{记号；组合分析。}

历史学和考古学使我们了解到为数众多的“记数体系”；它们最初的目标，是给每个单独的整数（直至一个视应用需要而定的界限）配一个名称和一种书面表示，后者由为数有限的符号按或多或少规则的法则组合而成。最为常见的做法是把整数分解为“相继单位”\(b_{1}, b_{2}, \ldots, b_{n}, \ldots\) 之和，其中每一个都是前一个的整数倍；虽然一般取 \(b_{n}/b_{n-1}\) 等于同一个数 b（体系的“基”，最常见的是 10），但这一规则有不少例外，如巴比伦人那里 \(b_{n}/b_{n-1}\) 有时等于 10，有时等于 6 {[}232{]}，又如玛雅人的历法体系中，\(b_{n}/b_{n-1}\) 除 n = 2 外都等于 20，而 \(b_{2}/b_{1} = 18\) {[}228{]}。至于相应的书面表示，它必须指明每一阶 i 的“单位”\(b_{i}\) 的数目；在许多体系中（如埃及人、希腊人和罗马人那里），相继的倍数 \(k.b_{i}\)（其中 k 从 1 变到 \((b_{i+1}/b_{i}) - 1\)）由一些同时依赖于 k 和 i 的符号来表示。第一个重要的改进在于：用同一个符号表示所有的数 \(k.b_{i}\)（对 k 的同一个值）；这就是“位值记数法”的原理，其中指标 i 由如下事实来指示：表示 \(k.b_{i}\) 的符号出现在构成所表示之数的诸“切片”序列中的第 i 个位置上。这种类型的第一个体系见于巴比伦人那里：无疑早在公元前 2000 年，他们就用同一个记号表示与指数 i 的一切值相应的所有倍数 \(k.60^{\pm i}\)（{[}232{]}，第~93-109~页）。这种体系的不便之处当然在于所用符号的歧义性：只要没有任何东西指明某一阶的单位可能空缺，换句话说，只要该体系未通过引入“零”而加以补全，歧义就一直存在。然而人们看到，巴比伦人在其历史的大部分时间里并不使用这样的符号；事实上，他们只是在公元前最后两个世纪才使用这样一种“零”，而且还只在数的内部使用；在那之前，只能靠上下文来确定所论之数的意义。只有另外两个体系系统地使用“零”：玛雅人的体系（据推测在公元纪年开始时即已在使用 {[}228{]}），以及我们现今的十进制体系，它（经由阿拉伯人）源自印度数学，在印度，其使用早在公元最初几个世纪即有明证。还应指出，把零当作一个数（而不只是一个简单的分隔记号）的概念及其引入运算，也属于印度人的独创性贡献（{[}78{]}，第 I 卷）。当然，一旦掌握了位值记数法原理，把它推广到任何基都是容易的事；关于 17 世纪以来提出的各种“基”孰优孰劣的讨论属于数值计算的技术问题，不宜在此展开。我们只限于指出，作为这些体系之基础的那个运算，即所谓“欧几里得”除法，在希腊人之前并未出现，它无疑应追溯到最早的一批毕达哥拉斯学派学者，他们把它作为其理论算术的基本工具（见第~85~页）。

冠以“组合分析”之名的那些一般的计数问题，看来直到古典古代的最后几个世纪才得到处理：只有公式 \(\binom{n}{2}=n(n-1)/2\) 在公元 3 世纪得到确证。印度数学家婆什迦罗（12 世纪）知道 \(\binom{n}{p}\) 的一般公式。更系统的研究见于 13 世纪初利维·本·热尔松的一部手稿：他得到了使 n 个对象中每次取 p 个的排列数 \(V_{n}^{p}\) 得以计算的递推公式，特别是 n 个对象的排列数；他还陈述了等价于关系 \(\binom{n}{p}=V_{n}^{p}/p!\) 与 \(\binom{n}{n-p}=\binom{n}{p}\) 的法则（{[}311{]}，第 VI 卷，第~64-65~页）。但这部手稿似乎被同时代人所忽视，其结果只是在后来的几个世纪中才被数学家们逐渐重新发现。在后起的进展中，我们可以指出：卡尔丹证明了 n 个元素的集合的非空子集数是 \(2^{n}-1\)；帕斯卡和费马在创立概率演算时重新发现了 \(\binom{n}{p}\) 的表达式，而帕斯卡是第一个注意到这些数与二项式公式之间关系的人：后者看来早在 13 世纪已为阿拉伯人所知，14 世纪为中国人所知，并于 16 世纪初在西方被重新发现，同时被重新发现的还有通常归于帕斯卡名下的、称为“算术三角形”的递推计算方法（{[}311{]}，第 VI 卷，第~35-38~页）。最后，莱布尼茨在 1676 年前后得到了（但未发表）“多项式系数”的一般公式，该公式由德·莫弗独立地重新发现，并于 20 年后发表。

